\begin{mistake}{2026-05-07}{04}

\begin{problemcontent}
设 \(x_1=\frac12\)，\(x_{n+1}=x_n^2+x_n\)（\(n=1,2,\cdots\)），求极限
\[
\lim_{n\to\infty}\left(\frac{1}{x_1+1}+\frac{1}{x_2+1}+\cdots+\frac{1}{x_n+1}\right).
\]
\end{problemcontent}

\begin{solutioncontent}
由递推式
\[
x_{k+1}=x_k^2+x_k=x_k(x_k+1),
\]
两边取倒数，得
\[
\frac{1}{x_{k+1}}=\frac{1}{x_k(x_k+1)}.
\]
利用部分分式分解
\[
\frac{1}{x_k(x_k+1)}=\frac{1}{x_k}-\frac{1}{x_k+1},
\]
所以
\[
\frac{1}{x_k+1}=\frac{1}{x_k}-\frac{1}{x_{k+1}}.
\]
于是前 \(n\) 项和为
\[
\begin{aligned}
S_n&=\sum_{k=1}^n\frac{1}{x_k+1}\\
&=\sum_{k=1}^n\left(\frac{1}{x_k}-\frac{1}{x_{k+1}}\right)\\
&=\frac{1}{x_1}-\frac{1}{x_{n+1}}.
\end{aligned}
\]
又因为 \(x_1=\frac12>0\)，且
\[
x_{n+1}-x_n=x_n^2>0,
\]
故 \(\{x_n\}\) 严格递增且 \(x_n\ge \frac12\)。若 \(\{x_n\}\) 有界，则单调有界必收敛，设极限为 \(A\)，由递推式取极限得
\[
A=A^2+A,
\]
从而 \(A=0\)，这与 \(x_n\ge \frac12\) 矛盾。因此 \(x_n\to +\infty\)，所以
\[
\frac{1}{x_{n+1}}\to 0.
\]
因此
\[
\lim_{n\to\infty}S_n=\frac{1}{x_1}-0=\frac{1}{1/2}=2.
\]
\end{solutioncontent}

\mistakeknowledge{
\begin{itemize}
  \item \textbf{核心公式/定理}：由 \(x_{k+1}=x_k(x_k+1)\) 得 \(\frac{1}{x_k+1}=\frac{1}{x_k}-\frac{1}{x_{k+1}}\)，从而裂项相消。
  \item \textbf{易错点}：不要直接估计 \(\frac{1}{x_k+1}\)；应先把目标项和递推式联系起来，对 \(x_{k+1}=x_k(x_k+1)\) 取倒数并部分分式分解。
  \item \textbf{关联知识}：证明 \(\frac{1}{x_{n+1}}\to0\) 时，可用 \(x_{n+1}-x_n=x_n^2>0\) 得单调递增，再用反证法排除有界情形，推出 \(x_n\to+\infty\)。
\end{itemize}}

\end{mistake}
